\chapter*{GLOSARIO}
\thispagestyle{fancy} 
\addcontentsline{toc}{chapter}{GLOSARIO}

\begin{longtable}{|p{0.2\textwidth}|p{0.75\textwidth}|}
	\hline
	\textbf{Término} & \textbf{Descripción} \\
	\hline
	\endfirsthead
	
	\hline
	\textbf{Término} & \textbf{Descripción} \\
	\hline
	\endhead
	
	\hline
	\endfoot
	
	3GPP & \textit{3rd Generation Partnership Project}. Organización de estandarización que desarrolla protocolos para telecomunicaciones móviles, incluyendo las especificaciones para redes 5G. \\
	\hline
	
	5G & Quinta generación de tecnología de redes móviles celulares, sucesora de 4G/LTE. \\
	\hline
	
	AMF & \textit{Access and Mobility Management Function}. Función de red del 5G Core que gestiona el acceso y movilidad de los equipos de usuario. \\
	\hline
	
	API & \textit{Application Programming Interface}. Interfaz de programación de aplicaciones que permite la comunicación entre diferentes componentes de software. \\
	\hline
	
	CNI & \textit{Container Network Interface}. Especificación y bibliotecas para configurar interfaces de red en contenedores Linux. \\
	\hline
	
	eMBB & \textit{Enhanced Mobile Broadband}. Tipo de \textit{slice} 5G optimizado para servicios de alto ancho de banda como video 4K/8K y realidad virtual. \\
	\hline
	
	gNodeB & Estación base de nueva generación en redes 5G, equivalente al eNodeB de LTE. \\
	\hline
	
	IoT & \textit{Internet of Things}. Internet de las Cosas, red de dispositivos físicos conectados que intercambian datos. \\
	\hline
	
	MANO & \textit{Management and Orchestration}. Marco de gestión y orquestación definido por ETSI para funciones de red virtualizadas. \\
	\hline
	
	mMTC & \textit{massive Machine-Type Communications}. Tipo de \textit{slice} 5G optimizado para conectar masivamente dispositivos IoT con bajo consumo energético. \\
	\hline
	
	NFV & \textit{Network Functions Virtualization}. Virtualización de funciones de red, permitiendo que servicios de red se ejecuten en software sobre hardware estándar. \\
	\hline
	
	NFVO & \textit{NFV Orchestrator}. Componente del marco MANO responsable de orquestar recursos de red y servicios. \\
	\hline
	
	NRF & \textit{Network Repository Function}. Función de red del 5G Core que mantiene un registro de funciones de red disponibles. \\
	\hline
	
	NSD & \textit{Network Service Descriptor}. Descriptor que define la estructura y requisitos de un servicio de red en formato YAML/TOSCA. \\
	\hline
	
	ONAP & \textit{Open Network Automation Platform}. Plataforma open-source de orquestación y automatización de redes de grado carrier. \\
	\hline
	
	OSM & \textit{Open Source MANO}. Plataforma open-source que implementa el marco ETSI NFV MANO para orquestación de funciones de red virtualizadas. \\
	\hline
	
	QoS & \textit{Quality of Service}. Calidad de Servicio, conjunto de tecnologías y mecanismos para garantizar niveles específicos de rendimiento de red. \\
	\hline
	
	RAN & \textit{Radio Access Network}. Red de Acceso Radio, componente de red móvil que conecta dispositivos de usuario mediante interfaces de radio. \\
	\hline
	
	SDR & \textit{Software-Defined Radio}. Radio definido por software, donde componentes tradicionalmente implementados en hardware se implementan mediante software. \\
	\hline
	
	SMF & \textit{Session Management Function}. Función de red del 5G Core responsable de gestionar sesiones de datos de usuario. \\
	\hline
	
	SST & \textit{Slice/Service Type}. Valor numérico que identifica el tipo de \textit{slice} de red según las categorías definidas por 3GPP (1=eMBB, 2=URLLC, 3=mMTC). \\
	\hline
	
	TOSCA & \textit{Topology and Orchestration Specification for Cloud Applications}. Lenguaje estándar para describir topologías y orquestación de aplicaciones en la nube. \\
	\hline
	
	UE & \textit{User Equipment}. Equipo de Usuario, terminal o dispositivo que se conecta a la red móvil. \\
	\hline
	
	UPF & \textit{User Plane Function}. Función de red del 5G Core que maneja el plano de usuario, incluyendo enrutamiento y reenvío de paquetes. \\
	\hline
	
	URLLC & \textit{Ultra-Reliable Low-Latency Communications}. Tipo de \textit{slice} 5G optimizado para aplicaciones críticas que requieren latencia ultrabaja y alta confiabilidad. \\
	\hline
	
	VIM & \textit{Virtualized Infrastructure Manager}. Gestor de infraestructura virtualizada que controla recursos de cómputo, almacenamiento y red. \\
	\hline
	
	VNFD & \textit{VNF Descriptor}. Descriptor que define las características y requisitos de una función de red virtualizada. \\
	\hline
	
	VNFM & \textit{VNF Manager}. Gestor de funciones de red virtualizadas, responsable del ciclo de vida de VNFs. \\
	\hline
	
\end{longtable}

\clearpage